\section{Introduction}\label{sec:introduction}
An optimization model can be numerically feasible and operationally unacceptable.
In a predictive controller, relaxing an actuator bound changes what the plant is
permitted to do; relaxing a preferred room temperature changes the service being
requested. Treating both departures as interchangeable slack conceals a policy
decision. The application studied here makes that decision explicit: physical
limits remain hard, while selected interval preferences receive specified lower-
and upper-side prices. The intended use is the construction and inspection of
small, reproducible optimization models before a practitioner embeds one in a
larger controller. The examples are synthetic. They do not measure occupant
satisfaction, building savings, industrial reliability or operational safety.

The numerical foundation is a convex quadratic program with sparse affine rows.
That foundation appears in constrained model predictive control, energy scheduling
and repeated allocation problems~\cite{rawlings,mayne,mayne}. A soft constraint
adds a convex price to selected departures rather than changing the definition of
physical feasibility. The important modeling questions are therefore which rows
may yield, which direction matters, which units a departure has, and which cost
is actually optimized. Exact penalty theory explains why an adequately large
linear price can sometimes recover a hard-constrained optimum; it does not say
that an arbitrary finite price makes every returned elastic optimum acceptable
for a physical system~\cite{kerrigan,nocedal}.

The implementation compares explicit lifting, generalized slack projection,
side-specific projection, DAQP, a fixed-parameter FlexQP reconstruction and
ElastiQP on their genuine compatibility domains. A separate directional group
peak implementation couples several interval departures while retaining the
original interval rows in its principal system. Neither unequal side prices nor
the ability to express a directional peak establishes a new modeling capability.
DAQP already supports unequal side weights; directional row splitting and a
compact epigraph express the grouped objective exactly~\cite{daqp,gsp,elasti}.
The computational question is narrower: does a particular original-row proximal
implementation reduce the cost of solving the same mathematical problem?

Two completed studies answer that question negatively for their registered
primary criteria. The first contains 456 seeded instances and retains 11592
accepted and 7560 incompatible measured observations. Its 168 accepted held-out
pairs yield a median relative cold end-to-end reduction of \OriginalEstimate\%,
with a 95\% seed-cluster interval from \OriginalLower\% to \OriginalUpper\%.
The second contains 210 repeated-forecast instances. Its complete 144-instance
held-out endpoint yields \PrimaryEstimate\%, with corresponding limits
\PrimaryLower\% and \PrimaryUpper\%, including \PrimaryFailed{} failed pairs.
Neither result establishes the predicted 10\% improvement. These are distinct
studies with different native packages, populations and timing boundaries.

The technical contribution is an inspectable implementation and evidence chain,
not a claim of uniform superiority. Specifically, the system provides an original-
space model with physical row metadata; exact adapters with explicit incompatibility
states; independent objective, hard-row and primal and dual checks; repeated-operation
measurements with retained failures; and versioned producer/consumer boundaries.
The distinction between numerical evidence and its presentation is implemented:
native code computes solutions, local tools export validated results, and the
static reader displays or imports matching evidence. A dormant API and unavailable
coupled browser path remain visibly unavailable.

The remainder gives the mathematical formulation and notation, details every
data exchange and processing stage, derives each executable method, and states
the experimental protocols. The results include model agreement, formulation
size, conditioning, repeated control and complete negative primary inference.
Validation and discussion distinguish a correct optimum from a fast computation
and from an operationally useful decision. The accompanying research paper
develops the two related proximal proposals and their separately attributed
experiments; this technical account also describes the full software system and
all of its contracts.
